\documentclass[10pt,a4paper]{amsart}
\usepackage[margin=19mm]{geometry}
\usepackage{amsmath,amssymb,mathtools}
\usepackage[scr=rsfs,cal=euler]{mathalfa}
\usepackage{xfrac}
\usepackage[unicode,hidelinks,bookmarks=false]{hyperref}
\newcommand{\ie}{i.e.\ }
\newcommand{\cf}{cf.\ }
\newcommand{\resp}{respectively\ }
\newcommand{\Subsection}{\subsection}
\providecommand{\nobreakdash}{\nobreak\hbox{-}}
\setlength{\emergencystretch}{2em}
\multlinegap0pt
\usepackage{bm}
\usepackage{bbm}
\usepackage{dsfont}
\usepackage{xspace}
\usepackage{cancel}
\usepackage{mathtools}
\usepackage{amsthm}
\newtheorem{theorem}{Theorem}
\newtheorem{lemma}[theorem]{Lemma}
\newcommand{\N}{\mathbb{N}}
\newcommand{\R}{\mathbb{R}}
\newcommand{\T}{\intercal}
\newcommand{\bx}{\bm{x}}
\newcommand{\la}{\lambda}
\newcommand{\sset}[1]{\left\{{#1}\right\}}
\newcommand{\tbx}{\tilde{\bm{x}}}
\title[Beyond the classification theorem of Cameron, Goethals, Seid]{Beyond the classification theorem of Cameron, Goethals, Seidel, and Shult}
\author{Hricha Acharya, Zilin Jiang}
\date{Source: arXiv:2404.13136v3 (CC BY 4.0)}
\begin{document}
\maketitle

\noindent\emph{Source and licence.} Beyond the classification theorem of Cameron, Goethals, Seidel, and Shult. Version arXiv:2404.13136v3 (CC BY 4.0), distributed under the Creative Commons Attribution 4.0 International licence, as recorded in the arXiv licence metadata for this exact version and shown on the abstract page https://arxiv.org/abs/2404.13136v3. Full rights notice: licenses/arxiv-2404.13136-CC-BY-4.0.txt.\par\smallskip

\noindent\textit{Editorial scope.} Counted unit: the Lemma whose source label is \texttt{lem:clique-extension} in the article (source0.tex lines 337--339, source label \texttt{lem:clique-extension}), delivered together with the article's own complete original proof of it (source0.tex lines 340--390, one \texttt{proof} environment of 51 lines). No mathematics is written, repaired or re-derived: statement and proof are the article's own text line for line. Required context retained uncounted: none. Every definition, display and label of those lines is retained unchanged. Omitted from this package and not redistributed: 1--336, 391--2881. Nothing inside a retained excerpt is omitted or altered: every retained body is delivered exactly as the article's lines are written, so no omission receipt of any kind accompanies this package. Citations: the retained excerpts carry no citation command at all, so no bibliography entry of the article is transcribed and no reference is redistributed. Reference adaptation: none was needed and none was made; every reference inside the delivered unit resolves inside it, so no unresolved reference or citation is produced. Printed slips kept literal: no printed word, symbol, hypothesis or number of the counted statement or of its proof was altered. Layout and class: the article's own class is \texttt{article} with packages including hyperref, amsfonts, amsmath, amssymb, amsthm, fullpage, enumitem, tikz, bm, mathtools, microtype, mleftright, cleveref; the wrapper is the lane's pinned \texttt{amsart} editorial wrapper at 10pt on A4, which restates the theorem-like environments the retained text needs and adds the article's own macro definitions that the retained text needs (\texttt{\textbackslash N}, \texttt{\textbackslash R}, \texttt{\textbackslash T}, \texttt{\textbackslash bx}, \texttt{\textbackslash la}, \texttt{\textbackslash sset}, \texttt{\textbackslash tbx}), with extra wrapper packages (bm, bbm, dsfont, xspace, cancel, mathtools). The delivered numbering is the unit's own. Assets: no retained span contains a figure environment, a graphics-inclusion command, a PDF-page inclusion, a table or a TikZ picture; no image file is copied into this package, the package declares no asset dependency and no image file is redistributed with it. Licence: version arXiv:2404.13136v3 of this article is distributed under the Creative Commons Attribution 4.0 International licence, as recorded in the arXiv licence metadata for this exact version and shown on the abstract page https://arxiv.org/abs/2404.13136v3. A full rights notice is delivered at licenses/arxiv-2404.13136-CC-BY-4.0.txt. Characters: every retained excerpt is pure ASCII, so the delivered engine, XeLaTeX with its default Latin Modern text font, reports no missing character.
\begin{lemma} \label{lem:clique-extension}
    For every rooted graph $F_R$ and $\la \ge 2$, if $\lim_{\ell\to\infty} \la_1(F_R, \ell) < -\la$, then there exists $m \in \N^+$ such that $\la_1(F_R, K_m) < -\la$.
\end{lemma}

\begin{proof}
    Pick $\ell \in \N^+$ such that $\la_1(F_R, \ell) < -\la$. Set $\la' := -\la_1(F_R, \ell)$. Let $v_0\dots v_\ell$ be the path added to $F$ to obtain $(F_R, \ell)$, where the vertex $v_0$ is connected to every vertex in $R$, and let $\bx\colon V(F) \cup \sset{v_0, \dots, v_\ell} \to \R$ be an eigenvector of $(F_R, \ell)$ associated with $-\la'$. We abuse notation and write $x_i$ in place of $x_{v_i}$ for $i \in \sset{0,\dots, \ell}$. Define $\tbx \colon V(F) \cup V(K_m) \to \R$ by
    \[
        \tilde{x}_v = \begin{cases}
            x_v & \text{if }v \in V(F); \\
            x_0/m  & \text{if }v \in V(K_m).
        \end{cases}
    \]

    We claim that $\sum_{v \in V(F)}x_v^2 > 0$. Indeed, assume for the sake of contradiction that $x_v = 0$ for $v \in V(F)$. Using $-\la' x_i = \sum_{u\sim v_i}x_u$ for $i \in \sset{0, \dots, \ell}$, where the sum is taken over all vertices $u$ that are adjacent to $v_i$ in $(F_R, \ell)$, we obtain that \[
        (A_P + \la' I)\begin{pmatrix}
            x_0 \\
            \vdots \\
            x_\ell
        \end{pmatrix} = \bm{0},
    \]
    where $P$ denotes the path $v_0\dots v_\ell$. Since $\la_1(P) > -2 > -\la'$, the matrix $A_P + \la' I$ is positive definite, contradicting the assumption that $\bx$ is nonzero.

    Because $\sum_{v \in V(F)}x_v^2 > 0$, clearly $\tbx$ is a nonzero vector. We compute
    \[
        \tbx^\T \tbx = \sum_{v \in V(F)}x_v^2 + m(x_0/m)^2.
    \]
    Moreover we can compute $\tbx^\T A_{(F_R, K_m)}\tbx$ as follows
    \[
        \tbx^\T A_{(F_R, K_m)}\tbx = \sum_{u,v \in V(F_R, 0)\colon u\sim v}x_ux_v + m(m-1)(x_0/m)^2.
    \]
    Since $\bx$ is an eigenvector of $(F_R, \ell)$ associated with $-\la'$, we obtain that
    \[
        \sum_{u, v \in V(F_R, 0)\colon u\sim v} x_ux_v + x_0x_1 = \sum_{v \in V(F_R, 0)}x_v \sum_{u\sim v}x_u = -\la'\sum_{v \in V(F)}x_v^2 - \la' x_0^2.
    \]
    Thus $\tbx^\T A_{(F_R, K_m)}\tbx$ can be simplified to
    \[
        \tbx^\T A_{(F_R, K_m)}\tbx = -\la'\sum_{v \in V(F)}x_v^2 -\la'x_0^2 - x_0x_1 + m(m-1)(x_0/m)^2.
    \]
    The Rayleigh principle says that $\la_1(F_R, K_m)$ is at most
    \[
        \frac{-\la'\sum_{v \in V(F)}x_v^2 - \la'x_0^2 - x_0x_1 + m(m-1)(x_0/m)^2}{\sum_{v \in V(F)}x_v^2 + m(x_0/m)^2},
    \]
    which, as $m\to\infty$, approaches
    \[
        -\la' - \frac{(\la'-2)x_0^2+x_0(x_0 + x_1)}{\sum_{v \in V(F)}x_v^2}.
    \]
    Here we used the claim that the denominator in the limit is positive.

    Recall that $\la' = -\la_1(F_R, \ell) > \la \ge 2$. It suffices to show that $x_0(x_0 + x_1)\ge 0$. In fact, we prove inductively that $x_i(x_i + x_{i+1}) \ge 0$ for $i \in \sset{\ell-1, \dots, 0}$. The base case where $i = \ell-1$ follows immediately from $-\la' x_\ell = x_{\ell-1}$ and $\la' > 2$. For the inductive step, using $-\la' x_{i+1} = x_i + x_{i+2}$ and $\la' > 2$, we obtain
    \begin{multline*}
        x_{i+1}(x_{i+1} + x_{i+2}) = x_{i+1}(x_{i+1} - \la'x_{i+1}-x_i) = -(\la'-2)x_{i+1}^2 - (x_i + x_{i+1})x_{i+1} \\
        \le -(x_i + x_{i+1})x_{i+1} = -(x_i + x_{i+1})^2 + x_i(x_i + x_{i+1}) \le x_i(x_i + x_{i+1}),
    \end{multline*}
    which implies that $x_i(x_i+x_{i+1}) \ge 0$ by the inductive hypothesis.
\end{proof}

\end{document}
